\section{Related Work}\label{sec:related}

\subsection{Classical network interdiction}\label{sec:rel-classical}

Maximum flow interdiction has been studied since the 1960s, first in military and logistics applications~\cite{ausiello2026}. The surveys of Smith and Song~\cite{smith2020} and of Ausiello \emph{et al.}~\cite{ausiello2026} cover the many variants proposed since then. Wood~\cite{wood1993} proposed the single-level MILP of Section~\ref{sec:bg-mfnip}, together with valid inequalities and the complexity results cited there. Phillips~\cite{phillips1993} studied the problem under the name network inhibition and gave pseudopolynomial algorithms for planar graphs.

In practice, exact methods rely on mathematical programming. Smith and Song~\cite{smith2020} identify models obtained by dualization, such as~\eqref{eq:milp}, and extended formulations solved by row generation as the most common approaches. Royset and Wood~\cite{royset2007} computed the efficient frontier between interdiction cost and maximum flow by Lagrangian relaxation combined with a specialized branch-and-bound that enumerates part of the $s$--$t$ cuts. Their method was one to two orders of magnitude faster than the same scheme with a standard integer programming solver.

\subsection{Quantum algorithms for graph problems}\label{sec:rel-graphs}

D\"urr, Heiligman, H{\o}yer, and Mhalla~\cite{durr2006} obtained nearly tight bounds on the quantum query complexity of connectivity, strong connectivity, minimum spanning tree, and single-source shortest paths, and their upper bounds use variants of quantum minimum finding. Ambainis and \v{S}palek~\cite{ambainis2006} presented quantum algorithms for matching and network flow. These works speed up the follower's problem, which is the maximum flow itself. Our approach is complementary, because the maximum flow remains a classical subroutine inside the oracle and we seek the quantum speedup in the interdictor's search over interdiction plans.

\subsection{Grover-based optimization}\label{sec:rel-gas}

Baritompa, Bulger, and Wood~\cite{baritompa2005} introduced Grover adaptive search (GAS), a framework that places the D\"urr--H{\o}yer method among threshold-based global optimization algorithms. Gilliam, Woerner, and Gonciulea~\cite{gilliam2021} applied GAS to constrained polynomial binary optimization, which includes QUBO problems~\cite{lucas2014,glover2019} as a special case. They built efficient oracles for polynomial objectives, placed the constraints inside the oracle, and tested the method on portfolio optimization in simulation and on quantum hardware. Grover search has also been applied to graph problems. Satoh, Ohkura, and Van Meter~\cite{satoh2020} solved a five-node maximum cut instance on IBM processors with a phase oracle divided into parts. In the GAS formulation of Gilliam \emph{et al.}, the search starts from the uniform superposition over all bit strings, and the oracle checks the constraints when there are any. Our interface differs in two respects. First, the oracle must solve a maximum flow problem, because the objective of the MFNIP is itself the value of an optimization problem. Second, the state preparation enforces the budget, so the search covers $\Fk$ instead of $\{0,1\}^{m}$.

\subsection{Quantum branch-and-bound}\label{sec:rel-limits}

A stronger rival is branch-and-bound, the method on which MILP solvers are built. Montanaro~\cite{montanaro2020} described a quantum algorithm that speeds up classical branch-and-bound almost quadratically in a fairly general setting. A speedup over the branch-and-bound tree explored when solving~\eqref{eq:milp} would be a result against a method used in practice rather than against enumeration. We leave this direction for future work (Section~\ref{sec:conclusion}).

\subsection{Search for prior work}\label{sec:rel-novelty}

To find out whether quantum algorithms had already been applied to network interdiction, we ran a systematic search on 6 November 2026 in four sources, namely the arXiv API (\texttt{export.arxiv.org}), the Semantic Scholar Graph API, the Crossref API, and a general web search engine. The queries combined the terms \emph{quantum}, \emph{Grover}, \emph{annealing}, or \emph{QAOA} with \emph{interdiction}, \emph{maximum flow}, \emph{bilevel}, \emph{most vital}, or \emph{edge removal}. For the arXiv we examined every result, and for the other sources the first 15 to 20 results by relevance. We considered relevant any work that combines a quantum algorithm with the removal of network elements. The exact queries and the number of results are listed in the repository (Section~\ref{sec:method-code}).

The search found two works that apply quantum algorithms to the removal of network elements. Peddireddy \emph{et al.}~\cite{peddireddy2026} formulated a bilevel vehicle routing problem with arc interdiction as a QUBO and solved it on D-Wave quantum annealers. Their problem concerns routing rather than flow, and their method is a heuristic without a performance guarantee. In a preprint, Sutcliffe and Mutyamsetty~\cite{sutcliffe2026} proposed selecting the connections to remove to contain the spread of malware by combining quantum amplitude estimation with Grover-based minimum finding over a candidate set of edge removals. Their method is the closest to ours, but the objective is a stochastic influence measure rather than a maximum flow, and no state preparation enforces the budget.

To the best of our knowledge, and based on this search, no prior work applies a quantum algorithm to the MFNIP, and none proposes a quantum search for interdiction in which the state preparation enforces the budget with proven query complexity. Because quantum methods for interdiction-type problems already exist~\cite{peddireddy2026,sutcliffe2026}, we restrict our novelty claim to the MFNIP and to the interface of Section~\ref{sec:interface}.
